%% Nomenclatura y notación. Va ANTES del índice y después de la leyenda de
%% niveles: PLAN-LIBRO.md §2.7 — nada se usa antes de estar definido.
\chapter*{Nomenclatura, niveles y notación}
\addcontentsline{toc}{chapter}{Nomenclatura, niveles y notación}

Este capítulo no supone nada. Da los cuatro alfabetos que hacen falta para leer el resto: las
\textbf{palabras}, los \textbf{niveles}, los \textbf{nombres} y los \textbf{signos}. Se puede
saltar y volver — pero nada de lo que viene después usa un término que no esté aquí, y eso está
comprobado por una máquina, no por buena voluntad.

\section*{1 · Las palabras}

Cada entrada lleva, además de su definición, una etiqueta que dice \textbf{de qué lado está lo que
nombra}. No es un adorno: es la primera pregunta que hay que saber contestar de cualquier palabra de
este libro, y dejarla implícita es la puerta por la que entró la inconsistencia que cuenta la
Parte~IV.

\begin{center}\footnotesize
\nivelterm{meta} del lado de Lean, donde razonamos\quad
\nivelterm{objeto} del lado del sistema aritmético del que se habla\\[0.35em]
\nivelterm{meta→objeto} una función de Lean que \emph{produce} sintaxis de ese sistema\quad
\nivelterm{ambos} nombra la distinción misma\\[0.35em]
\nivelterm{proyecto} una noción sobre el proyecto, no sobre las teorías
\end{center}

\input{extraido/_glosario.tex}

%% La leyenda de niveles. Va ANTES del índice: PLAN-LIBRO.md §2.5.
\section*{2 · Los dos ejes: quién afirma, y sobre qué}

Con esas palabras ya se puede plantear la pregunta que vuelve en cada línea:
\emph{¿esto es lenguaje objeto o lenguaje meta?} En un libro de metamatemática se
plantea en cada línea, y responderla mal es caro: la inconsistencia que cuenta la
Parte~IV fue exactamente un error de categoría. Por eso aquí el nivel \textbf{se imprime}.

No es una distinción binaria. Hay dos ejes independientes, y cada fragmento de código y cada
enunciado llevan una chapa con los dos:

\noindent\textbf{Eje 1 — ¿quién afirma?} Son cuatro relaciones de derivabilidad distintas, y
\emph{de tipos distintos}: lo que las separa es si arrastran contexto y si son recursivamente
enumerables.

\begin{center}\footnotesize\setlength{\tabcolsep}{4pt}
\begin{tabularx}{\linewidth}{@{}l >{\ttfamily\scriptsize}l cc >{\raggedright\arraybackslash}X@{}}
\toprule
chapa & tipo en Lean & ctx. & r.e. & papel \\
\midrule
\colorbox{nvLean}{\textcolor{white}{\sffamily\scriptsize\bfseries\strut\,Lean\,}} &
  \normalfont--- & --- & --- & la metateoría: un teorema \emph{sobre} el sistema \\
\colorbox{nvDer}{\textcolor{white}{\sffamily\scriptsize\bfseries\strut\,$\Gamma\vdash$\,}} &
  List Formula → Formula → Prop & sí & \textbf{no} & el cálculo $\omega$; sólido en $\mathbb{N}$ \\
\colorbox{nvPrf0}{\textcolor{white}{\sffamily\scriptsize\bfseries\strut\,Prf$_0$\,}} &
  Formula → Prop & no & sí & Hilbert \textbf{intuicionista} \\
\colorbox{nvPrf}{\textcolor{white}{\sffamily\scriptsize\bfseries\strut\,Prf\,}} &
  Formula → Prop & no & sí & Hilbert clásico: \textbf{el que se aritmetiza} \\
\colorbox{nvPrfH}{\textcolor{white}{\sffamily\scriptsize\bfseries\strut\,$\Gamma\,$PrfH\,}} &
  List Formula → Formula → Prop & sí & sí & variante con contexto; ahí vive la deducción \\
\bottomrule
\end{tabularx}
\end{center}

\noindent\textbf{Eje 2 — ¿sobre qué?} Cinco valores.

\begin{center}\small
\begin{tabularx}{\linewidth}{@{}l >{\raggedright\arraybackslash}X@{}}
\toprule
\colorbox{nvLean!18}{\textcolor{nvLean}{\sffamily\scriptsize\bfseries\strut\,Lean\,}} &
  objetos nativos: \ident{Nat}, listas, recursión estructural \\
\colorbox{nvLean!18}{\textcolor{nvLean}{\sffamily\scriptsize\bfseries\strut\,sintaxis\,}} &
  la representación en Lean de la sintaxis: \ident{Term}, \ident{Formula} \\
\colorbox{nvObj!18}{\textcolor{nvObj}{\sffamily\scriptsize\bfseries\strut\,objeto\,}} &
  enunciados de la teoría aritmética: $\obj{+}$, $\obj{\cdot}$, $\obj{\sigma}$, sus axiomas \\
\colorbox{nvCod!18}{\textcolor{nvCod}{\sffamily\scriptsize\bfseries\strut\,código\,}} &
  términos objeto que \emph{denotan} sintaxis: $\cod{\codigo{\varphi}}$, \ident{numeral} \\
\colorbox{nvProv!18}{\textcolor{nvProv}{\sffamily\scriptsize\bfseries\strut\,Prov\,}} &
  dentro del predicado de demostrabilidad: la teoría hablando de sí misma \\
\bottomrule
\end{tabularx}
\end{center}

\noindent \textbf{Los puentes, y el que no existe.} Las cuatro no son intercambiables, y la
dirección importa:
\[
  \Prf_0\,\varphi \;\longrightarrow\; \Prf\,\varphi \;\longrightarrow\;
  \bigl(\forall \Gamma,\ \mathsf{PrfH}\ \Gamma\ \varphi\bigr),
  \qquad
  \Prf\,\varphi \;\longrightarrow\; \text{\ident{axioms}} \vdash \varphi .
\]
El recíproco $\text{\ident{axioms}} \vdash \varphi \to \Prf\,\varphi$ \textbf{es falso, y tiene que
serlo}: si valiera, $\vdash$ sería r.e.\ y Tarski cerraría el paso. Ésa es la asimetría que obliga a
tener dos cálculos en vez de uno, y por eso la chapa distingue en cuál se afirma.
Además \ident{prf0_to_derives} no usa \emph{ninguna} meta-regla $\omega$ y \ident{prf_to_derives}
usa \ident{dne} exactamente una vez: la diferencia entre los dos puentes es, medida,
$\{\ident{FOL.MetaRules.dne}\}$.

\medskip\noindent Tres ejemplos que ocupan tres casillas distintas y que conviene comparar:
\ident{two_mul_consN} afirma \emph{en Lean} sobre \ident{Nat};
\ident{prf_formCode_numeral} afirma \emph{en Prf} sobre \textbf{códigos};
\ident{hFN} afirma lo mismo, pero \emph{en} $\vdash$.
Y \ident{pcc_dot_cons} afirma en \ident{Prf} \textbf{dentro de} \ident{Prov} — cuatro niveles de
distancia entre el primero y el último.

\begin{leccion}[Por qué la chapa es un control, no un adorno]
La cuarta fila del segundo eje es la que se cobró 31 módulos. \ident{tcFn} —«el código del
código»— es una operación sobre \textbf{sintaxis} que se declaró como función \textbf{objeto}, y
una función objeto sólo puede depender de \textbf{valores}. Con la chapa puesta, la mezcla se ve
antes de escribirla.
\end{leccion}


\section*{3 · Cómo se lee un nombre}

Los identificadores van en inglés y siguen las convenciones de Mathlib. Tres reglas bastan para
leer casi cualquiera:

\begin{enumerate}[leftmargin=1.4em]
\item \textbf{La conclusión primero, las hipótesis con \ident{_of_}.} El nombre dice \emph{qué se
  demuestra}, no cómo. \ident{isNat_succ_of_isNat} es «$\sigma n$ es natural, \emph{a partir de}
  que $n$ lo es».
\item \textbf{Los bicondicionales acaban en \ident{_iff}.} \ident{prf_In_iff_boundedIn} es un «si y
  sólo si».
\item \textbf{Los símbolos se deletrean}: \ident{add}, \ident{mul}, \ident{succ}, \ident{eq},
  \ident{lt}, \ident{le}, \ident{dvd}, \ident{not}, \ident{forall}, \ident{exists}. Así
  \ident{prf_add_zero_left} se lee sin traducir.
\end{enumerate}

Y hay algo más útil todavía, porque conecta con la chapa de la sección anterior:
\textbf{el prefijo dice en qué nivel se afirma.} No es una impresión — es la práctica medida del
proyecto:

\begin{center}\small
\begin{tabularx}{\linewidth}{@{}l r >{\raggedright\arraybackslash}X@{}}
\toprule
prefijo & símbolos & qué anuncia \\
\midrule
\ident{ax_}   & 154 & un \defterm{axioma objeto}: se postula, no se demuestra \\
\ident{prf_}  & 585 & un teorema del cálculo finitario \ident{Prf} \\
\ident{pcc_}  & 214 & un teorema de \ident{Prf} \textbf{dentro de} \ident{Prov} — la teoría hablando de sí misma \\
\ident{CRIT_} & 13  & un \defterm{sondeo} crítico: una pregunta contestada, a menudo en negativo \\
\bottomrule
\end{tabularx}
\end{center}

\nota{Los recuentos son de \modulo{ROBINSON\_PlusPlus/} a 2026-09-03. La excepción conocida son las
seis meta-reglas $\omega$ de \modulo{FOL}, anteriores a la convención, que no llevan \ident{ax_}.}

\section*{4 · Los signos}

\begin{center}\small
\begin{tabularx}{\linewidth}{@{}l >{\raggedright\arraybackslash}X@{}}
\toprule
signo & qué es \\
\midrule
$\Gamma \vdash A$ & $A$ se deriva de las hipótesis $\Gamma$ en el cálculo $\omega$ \\
$\Prf\,A$ & $A$ se deriva en el cálculo de Hilbert finitario, sin hipótesis \\
$\sigma n$ & el sucesor de $n$ \\
$t \doteq u$ & igualdad \textbf{del lenguaje objeto} (no la de Lean, que es \ident{=}) \\
\texttt{\#0} & la variable ligada más cercana: un \defterm{índice de De Bruijn} \\
$\codigo{\varphi}$ & el \defterm{código de Gödel} de $\varphi$: un término objeto que la denota \\
$\dotado{a}$ & «$a$ dotado»: el código de $a$ construido \emph{dentro} del lenguaje (\ident{tcFn}) \\
$\Prov(\codigo{\varphi})$ & la fórmula objeto que dice «$\varphi$ es demostrable» \\
$\Delta_0$, $\Sigma_1$ & clases de fórmulas por su cuantificación: acotada, y $\exists$ sobre acotada \\
\bottomrule
\end{tabularx}
\end{center}

\nota{Ningún término técnico aparece en este libro antes de estar definido aquí. Lo comprueba
\modulo{doc/book/scripts/terminos.py} en cada compilación, y no admite excepciones.}
